\documentclass[dvipdfmx,twoside]{jsarticle}
\usepackage{amsmath,amssymb}
\usepackage{CJKutf8}
\usepackage{okumacro}
\usepackage{booktabs}
\usepackage{array}
\usepackage{xcolor}
\usepackage{colortbl}
\usepackage{tipa}
\usepackage{tikz}
\usepackage{fancybox}
\usepackage{wasysym}
\usepackage{pifont}
\usepackage{graphicx}
\usepackage{wrapfig}
\usepackage{geometry}
\geometry{margin=2cm}
\usepackage{fancyhdr}
\usepackage[utf8]{inputenc}
\usepackage{CJK}
% 设置页面样式
\pagestyle{fancy}
\fancyhf{}
\renewcommand{\headrulewidth}{0pt}
\fancyhead[RO]{数学–\thepage}
\fancyhead[LE]{数学–\thepage}

% 重新定义 plain 样式
\fancypagestyle{plain}{%
  \fancyhf{}%
  \renewcommand{\headrulewidth}{0pt}%
  \fancyhead[RO]{数学-\thepage}%
  \fancyhead[LE]{数学-\thepage}%
}
% 自定义A方框命令
\newcommand{\abb}[1]{%
\begin{tikzpicture}[baseline]
\node[draw=black, 
      rectangle, 
      minimum width=0.8cm, 
      minimum height=0.3cm, 
      fill=gray!25, 
      font=\bfseries,
      line width=1pt,
      inner sep=2pt,
      anchor=base] {#1};
\end{tikzpicture}%
}
\newcommand{\ab}[1]{%
\begin{tikzpicture}[baseline]
\node[draw=black, 
      rectangle, 
      minimum width=0.8cm, 
      minimum height=0.3cm, 
      font=\bfseries,
      line width=1pt,
      inner sep=2pt,
      anchor=base] {#1};
\end{tikzpicture}%
}
\newcommand{\sab}[1]{%
\begin{tikzpicture}[baseline]
\node[draw=black, 
      rectangle, 
      minimum width=0.4cm, 
      minimum height=0.05cm, 
      font=\bfseries,
      line width=1pt,
      inner sep=0.5pt,
      anchor=base] {#1};
\end{tikzpicture}%
}
\newcommand{\maru}[1]{\tikz[baseline=-0.7ex]{
    \node[shape=circle,draw,inner sep=1pt,minimum size=5pt,anchor=center] {\footnotesize #1};}}
\definecolor{headercolor}{RGB}{220,220,220}
\definecolor{rowcolor1}{RGB}{245,245,245}
\definecolor{rowcolor2}{RGB}{255,255,255}

% \title{\vspace{-1.5cm} 2025年8月羚課文科数学月考 }
% \author{\textnosfal{Linc\ -\ 伊}}
\date{}
\begin{document}

\begin{CJK}{UTF8}{ipxm}  % 使用ipxm字体
  %月考表纸部分
\begin{center}

\vspace*{5cm}

% 学校logo
\includegraphics[width=5cm]{pics/1.jpg}

\vspace{2cm}

% 主标题
{\fontsize{24}{30}\selectfont\bfseries\sffamily
2025年9月\\
\vspace{1em}
羚課理科数学月考
}

\end{center}
\newpage
%第一部分
\noindent
\begin{tikzpicture}
\node[draw, rectangle, minimum width=1cm, minimum height=0.8cm, font=\Huge] {I};
\end{tikzpicture}
\\
\\
\textbf{問1}\qquad $a$, $b$ を定数とし、2次関数 $f(x) = x^2 + ax + b$ と1次関数 $h(x) = -3x$ を考える。関数 $y = h(x)$ のグラフと直線 $x = 3$ との交点を A、関数 $y = h(x)$ のグラフと直線 $x = -2$ との交点を B とする。関数 $y = f(x)$ のグラフの頂点 P が線分 AB 上にあるとき

$$L = \{f(3) - h(3)\} + \{f(-2) - h(-2)\}$$

の値の範囲を求めよう。

点 P の座標は

$$\left(-\frac{a}{\ab{\color{red}{\textsf{A}}}}, -\frac{a^2}{\ab{\color{red}{\textsf{B}}}} + b\right)$$

である。点 P が線分 AB 上にあるから、$a$ の値の範囲は

$$-\ab{\color{red}{\textsf{C}}} \leq a \leq \ab{\color{red}{\textsf{D}}}$$

である。また、$a$, $b$ は

$$b = \frac{a^2+\ab{\color{red}{\textsf{E}}} a}{\ab{\color{red}{\textsf{F}}}} $$

を満たす。よって、$L$ は $a$ を用いて

$$L = \frac{1}{2}a^2 + \ab{\color{red}{\textsf{G}}}a + \ab{\color{red}{\textsf{HI}}}$$

と表される。

したがって、$L$ の値の範囲は

$$\ab{\color{red}{\textsf{J}}} \leq L \leq \ab{\color{red}{\textsf{KL}}}$$
\newpage
\begin{center}
- 計算欄 (memo) -
\end{center}
\newpage
\noindent
\textbf{問2}\qquad 1から7までの数字が1つずつ書かれた7枚のカードが、左から小さい順に並んでいる。
この中から2枚のカードを選び、その位置を入れ換える操作を2回続けて行う。
2回の操作後の7枚のカードの並びを7桁の整数とみなすとき、この整数が偶数になる確率を求めよう。\\[1.5em]

\noindent
(1) \quad まず、1回目、2回目ともカードの入れ換えが7枚のカードの中から2枚を選べるときの確率を考える。\\[0.5em]

(i) \quad 1回目に7の書かれたカード以外のカードを入れ換える場合、2回の操作で偶数になる確率は $\frac{\ab{\color{red}{\textsf{M}}}}{\ab{\color{red}{\textsf{NO}}}}$ である。\\[0.5em]

(ii) \quad 1回目に7の書かれたカードを他のカードと入れ換える場合、2回の操作で偶数になる確率は $\frac{\ab{\color{red}{\textsf{P}}}}{\ab{\color{red}{\textsf{QR}}}}$ である。
\\[0.5em]

したがって、このとき2回の操作で偶数になる確率は $\frac{\ab{\color{red}{\textsf{ST}}}}{\ab{\color{red}{\textsf{UV}}}}$ である。\\[1.5em]

\noindent
(2) \quad 次に、2回目のカードの入れ換えでは、1回目に入れ換えた2枚を除いた残り7枚から2枚を選んで入れ換えるときの確率を考える。このとき2回の操作で偶数になる確率は $\frac{\ab{\color{red}{\textsf{W}}}}{\ab{\color{red}{\textsf{X}}}}$ である。
\newpage
\begin{center}
- 計算欄 (memo) -
\end{center}
\newpage
\noindent
\begin{tikzpicture}
\node[draw, rectangle, minimum width=1cm, minimum height=0.8cm, font=\Huge] {II};
\end{tikzpicture}
\\
\\
\textbf{問1}\qquad 2つの複素数 $\alpha, \beta$ は
\[
|\alpha| = 6, \quad |\beta| = 5, \quad |\alpha - \beta| = 6
\]
を満たすとする。\\[0.5em]

このとき
\[
\alpha \overline{\alpha} = \ab{\color{red}{\textsf{AB}}},\quad \beta \overline{\beta} = \ab{\color{red}{\textsf{CD}}},\quad \alpha \overline{\beta} + \overline{\alpha} \beta = \ab{\color{red}{\textsf{EF}}}
\]
を満たす。これらより
\[
|\alpha + \beta| = \sqrt{\ab{\color{red}{\textsf{GH}}}}
\]
となる。\\[0.5em]

ここで，複素数平面上の3点 $O(0)，A(\alpha)，B(\beta)$ を考えると
\[
\overrightarrow{OA}\cdot \overrightarrow{OB} = \frac{\ab{\color{red}{\textsf{IJ}}}}{\ab{\color{red}{\textsf{K}}}},\qquad \overrightarrow{AB} \cdot \overrightarrow{OB} = \frac{\ab{\color{red}{\textsf{LM}}}}{\ab{\color{red}{\textsf{N}}}}
\]
である。
\newpage
\begin{center}
- 計算欄 (memo) -
\end{center}
\newpage
\textbf{問2}\qquad 円 $C$ と直線 $\ell$, $m$ を

\begin{align*}
C &: x^2 + y^2 - 4x - 2y + 5 - r^2 = 0\\
\ell &: y = 2x - 2\\
m &: y = \frac{3}{4}x - 2
\end{align*}

とし、直線 $\ell$ に関して円 $C$ と対称な円を $C'$ とする。\\[1em]

(1) \quad 円 $C$ の中心は $\left(\ab{\color{red}{\textsf{O}}}, \ab{\color{red}{\textsf{P}}}\right)$、半径は $r$ である。\\[0.5em]

(2) \quad 円 $C'$ の中心 $(a,b)$ を求めよう。円 $C$ と円 $C'$ は直線 $\ell$ に関して対称であるから

$$a = \frac{\ab{\color{red}{\textsf{Q}}}}{\ab{\color{red}{\textsf{R}}}}, \quad b = \frac{\ab{\color{red}{\textsf{S}}}}{\ab{\color{red}{\textsf{T}}}}$$

である。\\[0.5em]

(3) \quad 円 $C'$ が直線 $m$ に接するように $r$ を求めよう。

$$r = \frac{\left|\ab{\color{red}{\textsf{U}}}a - \ab{\color{red}{\textsf{V}}}b + \ab{\color{red}{\textsf{W}}}\right|}{\ab{\color{red}{\textsf{X}}}}$$

であるから

$$r = \ab{\color{red}{\textsf{Y}}}$$

である。
\newpage
\begin{center}
- 計算欄 (memo) -
\end{center}
\newpage
\noindent
\begin{tikzpicture}
\node[draw, rectangle, minimum width=1cm, minimum height=0.8cm, font=\Huge] {III};
\end{tikzpicture}
\\
\\

$p > 1$, $q > 1$ とする。方程式

$$e^{2x} - ae^x + b = 0 \quad \cdots\cdots\cdots \maru{1}$$

において、$t = e^x$ とおくとき、$t$ に関する2次方程式

$$t^2 - at + b = 0$$

は解 $\log_{q^{2}} p$ と $\log_{p^{4}} q$ をもつとする。\\[0.5em]
\qquad このとき、$a$ の最小値とそのときの方程式①の解を求めよう。\\[2em]

(1) \quad まず

$$b = \dfrac{\ab{\color{red}{\textsf{A}}}}{\ab{\color{red}{\textsf{B}}}}$$

であり

$$a = \dfrac{\ab{\color{red}{\textsf{C}}}}{\ab{\color{red}{\textsf{D}}}} \log_q p + \dfrac{\ab{\color{red}{\textsf{E}}}}{\ab{\color{red}{\textsf{F}}}} \log_p q$$

である。\\[1em]

(2) \quad $p$, $q$ が $p > 1$, $q > 1$ を満たしながら動くとき、$\log_p q > \ab{\color{red}{\textsf{G}}}$ である。\\[0.5em]

したがって、$a$ は最小値 $\dfrac{\sqrt{\ab{\color{red}{\textsf{H}}}}}{\ab{\color{red}{\textsf{I}}}}$ を $\log_p q = \sqrt{\ab{\color{red}{\textsf{J}}}}$ のときにとる。\\[0.5em]

そのときの方程式①の解は

$$x = -\dfrac{\ab{\color{red}{\textsf{K}}}}{\ab{\color{red}{\textsf{L}}}} \log_e \ab{\color{red}{\textsf{M}}}$$

である。
\newpage
\begin{center}
- 計算欄 (memo) -
\end{center}
\newpage
\noindent
\begin{tikzpicture}
\node[draw, rectangle, minimum width=1cm, minimum height=0.8cm, font=\Huge] {IV};
\end{tikzpicture}
\\
\\
$$a_n = \int_0^1 x^{2n}\sqrt{1-x^2}\,dx \quad (n = 0, 1, 2, \cdots)$$

とおくとき、極限値 $\displaystyle\lim_{n \to \infty} \frac{a_n}{a_{n-1}}$ を求めよう。\\[0.5em]

(1) \quad まず、$a_0$, $a_1$ を求めてみよう。半径1の円の面積は $\pi$ であるから

$$a_0 = \int_0^1 \sqrt{1-x^2}\,dx = \frac{\pi}{\ab{\color{red}{\textsf{A}}}}$$

である。$a_1$ は部分積分法により
\begin{align*}
a_1 &= \int_0^1 x^2\sqrt{1-x^2}\,dx \\
  &= -\frac{\ab{\color{red}{\textsf{B}}}}{\ab{\color{red}{\textsf{C}}}}\left[x(1-x^2)^{\frac{\sab{\color{red}{\textsf{D}}}}{\sab{\color{red}{\textsf{E}}}}}\right]_0^1 + \frac{\ab{\color{red}{\textsf{F}}}}{\ab{\color{red}{\textsf{G}}}}\int_0^1 (1-x^2)^{\frac{\sab{\color{red}{\textsf{H}}}}{\sab{\color{red}{\textsf{I}}}}}\,dx \\
  &= \frac{\ab{\color{red}{\textsf{J}}}}{\ab{\color{red}{\textsf{K}}}}\int_0^1 \left(\sqrt{1-x^2} -  x^{\sab{\color{red}{\textsf{L}}}}\sqrt{1-x^2}\right)\,dx
\end{align*}

となる。よって、$a_1 = \frac{\pi}{\ab{\color{red}{\textsf{MN}}}}$ である。\\[1em]

(2) \quad 次の文中の $\ab{\color{red}{\textsf{O}}}\sim \ab{\color{red}{\textsf{U}}}$ には、下の選択肢 \maru{0} $\sim$ \maru{9} の中から適するものを選びなさい。\\[0.5em]

\qquad $a_1$ を求めたのと同様にして、$a_n$ は部分積分法により

$$a_n = \frac{\ab{\color{red}{\textsf{O}}}}{\ab{\color{red}{\textsf{P}}}}\int_0^1 \left( x^{\sab{\color{red}{\textsf{Q }}}}\sqrt{1-x^2} -  x^{\sab{\color{red}{\textsf{R}}}}\sqrt{1-x^2}\right)\,dx \quad (n = 1, 2, 3, \cdots)$$

となる。よって

$$\dfrac{a_n}{a_{n-1}} = \dfrac{\ab{\color{red}{\textsf{S}}}}{\ab{\color{red}{\textsf{T}}}}$$

となる。したがって

$$\lim_{n \to \infty} \frac{a_n}{a_{n-1}} = \ab{\color{red}{\textsf{U}}}$$

を得る。\\[0.5em]

\maru{0} \  $0$\quad
\maru{1} \  $1$\quad
\maru{2} \  $2$\quad
\maru{3} \  $3$\quad
\maru{4} \  $4$\quad
\\[0.5em]
\indent
\maru{5} \  $2n - 2$\quad
\maru{6} \  $2n - 1$\quad
\maru{7} \  $2n$\quad
\maru{8} \  $2n + 1$\quad
\maru{9} \  $2n + 2$\quad
\newpage
\begin{center}
- 計算欄 (memo) -
\end{center}

\end{CJK}
\end{document}